\section{引言：为什么评估会变成一门大课}

这节课讨论 \textbf{evaluation}。乍看之下，评估似乎很简单：给定一个固定模型，问它``有多好''，然后算一个分数就行了。实际情况远没有这么直接。评估不仅决定了论文表格里的数字，也在很大程度上决定了模型开发、产品选择、政策判断和安全测试的方向。

\begin{importantbox}{本课的中心观点}
评估不是一个机械步骤，而是一个需要先定义目标、再定义输入、调用方式、输出度量与结果解释的完整系统设计问题。
\end{importantbox}

课程一开始展示了几个很常见但含义并不完全相同的榜单：模型在 MMLU、MATH、GPQA 等基准上的分数，HELM 的能力榜，Artificial Analysis 的能力-价格前沿，OpenRouter 的流量排行，以及 Chatbot Arena 的人类偏好排名。它们都在回答``模型好不好''，但回答角度完全不同。一个数字本身没有意义，只有它对应的评价目标才有意义。

\subsection{你平时看到的评估长什么样}

现实中的评估通常包含这些元素：

\begin{itemize}
    \item \textbf{能力基准}：例如 MMLU、MATH、GPQA、Humanity's Last Exam。
    \item \textbf{偏好排序}：例如 Chatbot Arena 的 pairwise 比较与 ELO。
    \item \textbf{成本视角}：同样的准确率下，推理价格和训练成本可能差很多。
    \item \textbf{使用者视角}：模型不一定要``更聪明''，但一定要``更能卖得出去''或``更适合这个场景''。
    \item \textbf{安全视角}：拒答、越狱、幻觉、双重用途都属于评估范围。
\end{itemize}

\begin{knowledgebox}{评估并不只看准确率}
如果一个模型得分更高，但推理成本高 10 倍、部署风险更大、或者只是在训练数据上记住了测试集，那么它未必是更好的选择。评估必须把成本、鲁棒性和真实使用场景一起考虑进去。
\end{knowledgebox}

